\section{Procedimientos para análisis de alimentos}

\subsection{Proceso de Kjendalh}

\begin{enumerate}
    \item Adición de la muestra
    \item Digestión ácida con ácido sulfúrico ($H_{2}SO_{4}$) y algún catalizador (por ejemplo, $ZnCl$ o $BaCl$). Aquí se degrada la materia orgánica y sale en forma de gas, excepto por el nitrógeno ($N$), el cual permanece en forma de amoniaco ($NH_{3}$).
    \item Enfriamiento de la muestra
    \item Destilación\footnote{PENDIENTE DE REVISIÓN}
    \item Calibración de equipo (pHmetro)
    \item Medición de nitrógeno ($N$). Se realiza mediante una neutrilzación ácido-base con hidróxido de sodio ($NaOH$).
    \item Cálculo de nitrógeno perteneciente a proteína usando un factor dado. El factor estandar es 6.25
    \item Comparación con etiquetado\footnote{PENDIENTE DE REVISIÓN}
\end{enumerate}

\subsection{Determinación de lípidos}

\begin{enumerate}
    \item Se separan los lípidos de la matriz alimentaria usando un solvente dado (por ejemplo, éter)
    \item Se evapora el solvente
    \item Se pesa la materia restante
\end{enumerate}